%\documentclass[12pt,oneside]{amsart}
\documentclass{scrartcl}

%\documentclass{article}

\usepackage{amsthm,amsmath,amssymb}
\usepackage[numbers]{natbib}

%\usepackage[colorlinks]{hyperref}
\usepackage{hyperref}
\hypersetup{
    colorlinks,%
    citecolor=black,%
    filecolor=black,%
    linkcolor=black,%
    urlcolor=black
}
\usepackage{hypernat}
\usepackage[top=2.5cm, bottom=2.5cm, left=2.8cm,
right=2.8cm]{geometry} 

%\setlength{\parskip}{0pt}
%\setlength{\parsep}{0pt}
%\setlength{\headsep}{0pt}
%\setlength{\topskip}{0pt}
%\setlength{\topmargin}{0pt}
%\setlength{\topsep}{0pt}
%\setlength{\partopsep}{0pt}

\linespread{1}

%\usepackage{marvosym}

%\textwidth = 6.5 in \textheight = 9.2 in \oddsidemargin = 0.0 in
%\evensidemargin = 0.0 in \topmargin = -0.0 in \headheight = 0.0 in
%\headsep = 0.0in \parskip = 0.0in \parindent = 0.0in
%\def\baselinestretch{0.871}


\newtheorem{theorem}{Theorem}[section]
\newtheorem{proposition}[theorem]{Proposition}
\newtheorem{problem}[theorem]{Problem}
\newtheorem{fact}[theorem]{Fact}
\newtheorem{lemma}[theorem]{Lemma}
\newtheorem{defn}[theorem]{Definition}
\newtheorem{corollary}[theorem]{Corollary}
\newtheorem{remark}{Remark}[section]
\newtheorem{example}[theorem]{Example}

\newcommand{\E}{{\mathbb E}}
\newcommand{\se}{{\mathcal{E}}}
\newcommand{\R}{{\mathbb R}}
\renewcommand{\P}{{\mathbb P}}
\newcommand{\ac}{{\mathcal{A}}}
\newcommand{\A}{{\mathcal{A}}}
\newcommand{\B}{{\mathcal{B}}}
\newcommand{\C}{{\mathcal{C}}}
\newcommand{\M}{{\mathcal{M}}}
\newcommand{\Mf}{{\mathfrak{M}}}
\newcommand{\T}{{\mathcal{T}}}
\newcommand{\Z}{{\mathcal{Z}}}
\newcommand{\K}{{\mathcal{K}}}
\newcommand{\lc}{{\mathcal{L}}}
\newcommand{\Ps}{{\mathcal{P}}}
\newcommand{\G}{{\mathcal{G}}}
\newcommand{\pa}{{\mathring{p}}}
\newcommand{\F}{{\cal F}}
\newcommand{\samp}{{\mathcal{X}}}
\newcommand{\X}{{\mathcal{X}}}
\newcommand{\hi}{{\hat{\Phi}}}
\newcommand{\data}{{\text{data}}}

\newcounter{rcnt}[section]
\renewcommand{\thercnt}{(\roman{rcnt})}

\def\qt#1{\qquad\text{#1}}

\def\argmin{\mathop{\rm argmin}}
\def\argmax{\mathop{\rm argmax}}


\setlength{\parskip}{1.4 \medskipamount}
%\sloppy
%\linespread{1.3}

\begin{document}

\title{STAT 238 - Bayesian Statistics \\
 Homework One} 
\subtitle{Spring 2026, UC Berkeley}
\author{Due by 11:59 pm on 13 February 2026 \\ Total Points = 75} 

\date{}
\maketitle



% \tableofcontents

\begin{enumerate}

\item    Consider a clinical test for cancer that can yield either a
  positive (+)  
  or negative (-) result. Suppose that a patient who truly has cancer
  has a 1\% chance of slipping past the test undetected. On the other
  hand, suppose that a cancer-free patient has a 5\% probabiliity of
  getting a positive test result. Suppose also that 2\% of the
  population has cancer. Suppose that a patient who has taken the test
  \textbf{twice} has tested positive both times. 
  \begin{enumerate}
  \item What is the probability that this patient has cancer?
    (\textbf{3 points}). 
  \item What is the probability that this patient will test positive
    if they take the test the third time? (\textbf{3 points})
  \end{enumerate}
  Describe clearly the probability model you are using to
  calculate the above probabilities.

\item We have an urn with $N$ balls of which $\theta$ are red and the
  remaining $N - \theta$ are white. We are assuming that $N$ is known
  but $\theta$ is unknown. We sample $n$ balls from the urn without
  replacement and suppose we see $x$ red balls in the sample. What can
  we infer from $\theta$ based on knowledge of $x$, $n$ and $N$?

  This problem has applications in survey sampling. Suppose we want to
  know the number $\theta$ in a city (of total population $N$) that
  support a particular political party. We take a sample of $n$ people
  out of which $x$ support that party. What can we then say about
  $\theta$?

  \begin{enumerate}
  \item Argue that the likelihood function is given by: (\textbf{3 points})
    \begin{align*}
      \P \left\{X = x \mid N, \theta, n \right\} = \frac{{\theta
      \choose x}{N - \theta \choose n-x}}{{N \choose n}}
    \end{align*}
  This is the Hypergeometric distribution
  (\url{https://en.wikipedia.org/wiki/Hypergeometric_distribution}).  
  \item Work with the prior:
    \begin{align*}
      \P \left\{\theta = \theta \mid N \right\} = \frac{1}{N+1}
      I\{\theta \in \{0, 1, \dots, N\}\}, 
    \end{align*}
    and show that the posterior is given by (\textbf{3 points})
    \begin{align*}
      \P \left\{\theta = \theta \mid x, N, n \right\} = \frac{{\theta
      \choose x} {N - \theta \choose n - x}}{{N+1
      \choose n+1}} I\{\theta \in \{0, 1, \dots, N\}\}. 
    \end{align*}
  \item Show that (\textbf{3 points})
    \begin{align*}
      \E \left(\theta \mid x, n, N \right) = \frac{(N+2)(x+1)}{n+2} - 1
    \end{align*}
    and
    \begin{align*}
      \text{var}(\theta \mid x, n, N) = \frac{p(1 - p)}{n+3}(N+2)(N-n)
      \qt{where $p := \frac{x + 1}{n+2}$}. 
    \end{align*}
  \item When $N$ is large, argue that (\textbf{2 points})
    \begin{align*}
      \E \left(\frac{\theta}{N} \mid x, n, N \right) \approx
      \frac{x+1}{n+2} = p
    \end{align*}
    and
    \begin{align*}
      \text{var}\left(\frac{\theta}{N} \mid x, n, N \right) \approx
      \frac{p(1-p)}{n+3}. 
    \end{align*}
    This implies that, when $N$ is large, a point estimate of $\theta/N$ can be taken to
    be $p$ and the uncertainty of this point estimate can be taken to
    be the posterior standard deviation $\sqrt{\frac{p(1-p)}{n+3}}$.
  \item Let $R_{n+1}$ denote the event that the next draw (i.e., the
    $(n+1)$-th draw) leads to a red ball. Show that the posterior
    probability of $R_{n+1}$ given $x, n, N$ is given by: (\textbf{3 points})
    \begin{align*}
      \P \left\{R_{n+1} \mid x, n, N \right\} = \frac{x+1}{n+2} = p. 
    \end{align*}
    This formula is known as the Laplace Rule of Succession. 
  \end{enumerate}

\item In Lectures 3 and 4, we discussed the problem of estimating
  $\theta$ from numerical measurements $x_1, \dots, x_n$. We used the
  following Bayesian model:
  \begin{align*}
    X_1, \dots, X_n \mid \theta, \sigma \overset{\text{i.i.d}}{\sim}
    N(\theta, \sigma^2)  ~~ \text{ and } ~~ \theta, \log \sigma
    \overset{\text{i.i.d}}{\sim} \text{uniform}(-\infty, \infty).
  \end{align*}
  We proved that the posterior density of $\theta$ under this model
  satisfies:
  \begin{equation}\label{t_class}
    \frac{\sqrt{n} \left(\theta - \bar{x} \right)}{\hat{\sigma}} \mid
    \text{data} \sim t_{n-1}
  \end{equation}
  where
  \begin{equation*}
    \bar{x} := \frac{x_1 + \dots + x_n}{n} ~~ \text{ and } ~~
    \hat{\sigma}^2 := \frac{1}{n-1} \sum_{i=1}^n (x_i - \bar{x})^2. 
  \end{equation*}
  In this exercise, we shall go over a slightly different proof of the
  same result. Along the way, we shall also discuss inference for
  $\sigma$.
  \begin{enumerate}
  \item Show that
    \begin{align*}
      \theta \mid \text{data}, \sigma \sim N(\bar{x},
      \frac{\sigma^2}{n}). 
    \end{align*}
    The left hand side above refers to the conditional distribution of
    $\theta$ given the data as well as the parameter
    $\sigma$. (\textbf{2 points}).
  \item Show that the posterior density of $\sigma$ satisfies:
    (\textbf{2 points})
    \begin{align*}
   f_{\sigma \mid \text{data}}(\sigma) \propto   \sigma^{-n} \exp \left(-\frac{\sum_{i=1}^n (x_i -
      \bar{x})^2}{2 \sigma^2} \right) I\{\sigma > 0\}. 
    \end{align*}
  \item Use the results of the previous two parts along with the
    formula:
    \begin{align*}
      f_{\theta \mid \text{data}}(\theta) = \int_0^{\infty} f_{\theta
      \mid \text{data}, \sigma}(\theta) f_{\sigma \mid
      \text{data}}(\sigma) d\sigma
    \end{align*}
    to deduce \eqref{t_class}. (\textbf{2 points})
  \item Use the result in part (b) to show that
    \begin{align}\label{chi_square}
      \frac{1}{\sigma^2} \sum_{i=1}^n(x_i - \bar{x})^2 \mid
      \text{data} \sim \chi^2_{n-1}. 
    \end{align}
  On this basis, comment on whether $\hat{\sigma} = \sqrt{\sum_{i=1}^n (x_i -
      \bar{x})^2/(n-1)}$ is a reasonable point estimate of $\sigma$
    (\textbf{2 + 1 = 3 points}).
  \item Using \eqref{chi_square}, derive a $100(1-\alpha)\%$ Bayesian
    credible interval for $\sigma$ for fixed $\alpha$. (\textbf{2
      points}).
  \item Suppose $n = 6$ and $x_1 = 26.6, x_2 = 38.5, x_3 = 34.4,
  x_4 = 34, x_5 = 31, x_6 = 23.6$. Explicitly calculate the $95\%$
  credible interval for $\sigma$. (\textbf{2 points}). 
  \end{enumerate}

\item In Lecture 4, we considered the frequentist probability: 
  \begin{align}\label{freq_prob}
    \P\left\{\bar{X}_N -     \frac{t_{N-1,
    \alpha/2}}{\sqrt{N}}\sqrt{\frac{1}{N-1} \sum_{i=1}^N(X_i -
    \bar{X})^2} \leq \theta \leq \bar{X}_N +     \frac{t_{N-1,
    \alpha/2}}{\sqrt{N}}\sqrt{\frac{1}{N-1} \sum_{i=1}^N(X_i -
    \bar{X})^2} 
\right\}
  \end{align}
  where $X_1, X_2, \dots$ are i.i.d $N(\theta, \sigma^2)$ (for fixed
  $\theta$ and $\sigma$), and $N$ is the stopping rule:
  \begin{align}\label{Nstop}
    N := \inf\{n \geq 1 : X_n \leq 25\}. 
  \end{align}
  It was mentioned in class that, unlike the case when $N$ is a fixed number such as
  6, the probability \eqref{freq_prob} will be different from $1 -
  \alpha$. The goal of this problem is to verify this by simulation.

The probability \eqref{freq_prob} is actually not well-defined
 because of the presence of $N-1$ in the denominator, as $N$ in
 \eqref{Nstop} can equal 1 with some positive probability. We
 therefore attempt to calculate the simpler probability:
  \begin{align}\label{freq_prob_simp}
    \P\left\{\bar{X}_N -   \sigma  \frac{t_{N-1,
    \alpha/2}}{\sqrt{N}} \leq \theta \leq \bar{X}_N + \sigma    \frac{t_{N-1,
    \alpha/2}}{\sqrt{N}}  \right\}. 
  \end{align}
Follow the steps below to write code to calculate the above
probability.  
\begin{enumerate}
\item Fix $\theta = 35.5$ and $\sigma = 5.5$.
\item Set $\alpha = 0.05$ and $M = 100000$.
\item For $i = 1, \dots, M$:
  \begin{itemize}
  \item Generate $N(\theta,\sigma^2)$ random variables $X_1,
    X_2,\dots$ until an observation less than $25$ is obtained. Let
    $N$ be the number of observations.
  \item Calculate  $C_i = I\left\{\bar{X}_N -   \sigma  \frac{t_{N-1,
    \alpha/2}}{\sqrt{N}} \leq \theta \leq \bar{X}_N + \sigma    \frac{t_{N-1,
    \alpha/2}}{\sqrt{N}} \right\}$. 
\end{itemize}
\item The estimate for the probability \eqref{freq_prob_simp} is $(C_1
  + \dots + C_M)/M$. 
\end{enumerate}
Write code and compute this estimate of
\eqref{freq_prob_simp}. Comment on how different the estimate is from
$1 - \alpha$. (\textbf{5 points}).

\item Consider again the problem of estimating an unknown physical
  quantity $\theta$ from data $x_1 = 26.6, x_2 = 38.5, x_3 = 34.4,
  x_4 = 34, x_5 = 31, x_6 = 23.6$. Suppose that the measurement device
  can only record observations if they are in the range $[23,
  39]$. Based on this information and the data, what can be inferred
  about $\theta$?
\begin{enumerate}
\item Argue that the likelihood now should be: 
  \begin{align*}
    \prod_{i=1}^n \frac{\frac{1}{\sigma}
    \phi\left(\frac{x_i-\theta}{\sigma} \right)}{\Phi\left(\frac{39 -
    \theta}{\sigma} \right) - \Phi\left(\frac{23 -
    \theta}{\sigma} \right)}
  \end{align*}
  where $\phi(z) := (2 \pi)^{-1/2}\exp(-z^2/2)$ is the normal density
  and $\Phi(z) = \int_{-\infty}^z \phi(u) du$ is the normal
  cdf. (\textbf{2 points})
\item Using the same prior $\theta, \log \sigma
  \overset{\text{i.i.d}}{\sim} \text{uniform}(-\infty, \infty)$, write
  down the joint posterior density for $\theta$ and
  $\sigma$. (\textbf{2 points})
\item Construct a dense grid of points $(\theta, \log \sigma)$ (say in
  the range $[15, 50] \times [-1, 3]$). On this grid, compute the
  unnormalized posterior density, and normalize
  numerically. (\textbf{3 points})
\item Based on the discretized posterior, compute an approximate
  $95\%$ Bayesian credible interval for $\theta$. Compare your
  interval to the interval obtained by the usual normal analysis from
  class which does not take the truncation into consideration. For
  example, does this analysis with the truncated prior lead to a
  smaller credible region? (\textbf{4 points}). 
\end{enumerate}

\item Consider again the problem of estimating an unknown physical 
  quantity $\theta$ from data. Suppose the measurement device works
  correctly if the measurement lies in the range $[23, 39]$. If the
  measurement  falls outside this range, then the device displays
  \texttt{Failure}. Suppose we got the observations: $x_1 = 26.6, x_2 = 38.5, x_3 = 34.4,
  x_4 = 34, x_5 = 31, x_6 = 23.6, x_7 = \texttt{Failure}, x_8 =
  \texttt{Failure}$. Based on this information and the data, what can
  be inferred   about $\theta$?
\begin{enumerate}
\item Argue that the likelihood now should be: 
  \begin{align*}
\left[    \prod_{i=1}^6\frac{1}{\sigma}
    \phi\left(\frac{x_i-\theta}{\sigma} \right) \right] \left[1 -
    \Phi\left(\frac{39 - \theta}{\sigma} \right) +  \Phi\left(\frac{23
    - \theta}{\sigma} \right) \right]^2
  \end{align*}
  where $\phi(z) := (2 \pi)^{-1/2}\exp(-z^2/2)$ is the normal density
  and $\Phi(z) = \int_{-\infty}^z \phi(u) du$ is the normal
  cdf. (\textbf{2 points})
\item Using the same prior $\theta, \log \sigma
  \overset{\text{i.i.d}}{\sim} \text{uniform}(-\infty, \infty)$, write
  down the joint posterior density for $\theta$ and
  $\sigma$. (\textbf{2 points})
\item Construct a dense grid of points $(\theta, \log \sigma)$ (say in
  the range $[15, 50] \times [-1, 3]$). On this grid, compute the
  unnormalized posterior density, and normalize
  numerically. (\textbf{3 points})
\item Based on the discretized posterior, compute an approximate
  $95\%$ Bayesian credible interval for $\theta$. Compare your
  interval to the interval obtained by the usual normal analysis from
  class which does not take the truncation into consideration. For
  example, does this analysis with the truncated prior lead to a
  smaller credible region? (\textbf{4 points}). 
\end{enumerate}


  

\item Consider yet again the problem of estimating an unknown physical
  quantity $\theta$ from observations $x_1, \dots, x_n$. Now consider the
  following Bayesian model:
    \begin{align*}
    X_1, \dots, X_n \mid \theta, \sigma \overset{\text{i.i.d}}{\sim}
    \text{Lap}(\theta, \sigma)  ~~ \text{ and } ~~ \theta, \log \sigma
    \overset{\text{i.i.d}}{\sim} \text{uniform}(-\infty, \infty).
    \end{align*}
    Here $\text{Lap}(\theta, \sigma)$ denotes the Laplace distribution which has density
    \begin{equation*}
x \mapsto    \frac{1}{2\sigma} \exp \left(-\frac{|x - \theta|}{\sigma} \right).
    \end{equation*}
    \begin{enumerate}
    \item Show that the posterior of $\theta$ satisfies: (\textbf{4 points})
  \begin{equation*}
     f_{\theta \mid \text{data}}(\theta) \propto
     \left(\frac{1}{M(\theta)} \right)^n \qt{where $M(\theta) :=
       \sum_{i=1}^n |x_i - \theta|$}. 
   \end{equation*}
  \item For the data  $x_1 = 26.6, x_2 = 38.5, x_3 = 34.4,
  x_4 = 34, x_5 = 31, x_6 = 23.6$, evaluate the density at a dense
  grid of points in some large range (e.g., $[15, 50]$), normalize
  correctly so that the integral of the posterior density is 1, and
  then plot the posterior density. Comment on the nature of the
  plot. (\textbf{4 points}).
  \item Use your discretized approximation of the posterior to
    construct an approximate $95\%$ credible interval for
    $\theta$. Compare your interval with the interval derived from the
    normal based method from class. (\textbf{4 points}).
  \item Repeat the exercise from part (c) when the data is $x_1 =
    26.6, x_2 = 38.5, x_3 = 34.4, x_4 = 34, x_5 = 31, x_6 = 23.6, x_7
    = 120$. The additional data point $x_7 = 120$ is an outlier, and
    the goal here is to see how the outlier affects the two intervals based on the
    normal and  Laplace models. (\textbf{3 points}).
  \end{enumerate}

\end{enumerate}


\end{document}

%%% Local Variables:
%%% mode: latex
%%% TeX-master: t
%%% End:
